% planning/y7-q1/y7-q1.tex
%
% Year 7 Mathematics & Science -- lesson plan, Quarter 1 (from 4 August 2026).
% One week block is appended each Friday; Week 6 is the latest.
% Build with:  pwsh .claude/skills/weekly-plan/scripts/build-plan.ps1 y7-q1
%
% Structure, and it is the same for every week:
%   week overview page   -- at a glance, plus every live link for the week
%   parent update page   -- the Friday post, as it went out
%   one page per weekday -- the four-column table, one day per page, always
%
% Weeks 1 and 2 were transcribed from the Word plan; every week since was
% written here first, from the lessons' own plan.js files. Keep the ASCII rule: no literal
% en dashes, curly quotes, minus signs or times signs in this file. Use --,
% `` '', \tminus and \ttimes. See ../plan-style.tex.

\documentclass[11pt]{article}
% planning/plan-style.tex
% Shared house style for the printed weekly lesson-plan document. \input this
% from the plan file, right after \documentclass.
%
% Do not fork this file per document. If a plan needs a different look, the
% question is whether every plan should have it.
%
% The choices here are deliberate:
%   * LANDSCAPE A4. The day table is four columns of prose; portrait crushes
%     the Main Activities column into a ribbon and the page stops being usable
%     at a glance, which is the only thing this document is for.
%   * ONE DAY PER PAGE, always. A day that spills onto a second page is worse
%     than a day set two points smaller, so the table font is deliberately
%     small and the day page ends with \clearpage rather than \newpage.
%   * NO MATH MODE ANYWHERE. This document never loads mathastext, so a single
%     "$-3$" would render in Computer Modern serif against Lato and look
%     half-typeset. Use \tminus, \ttimes and \textsuperscript instead.
%   * Palette and link colours match the lesson decks and the homework packets.
%
% Provides: \tminus \ttimes \sep \playmark, the \daypage environment with
% \subjectrow, \weekoverview \glancerow \linkline, \announcepage \annsec
% \annsub, the hand-rolled contents macros \tocweek \tocline, and \slidelink
% \planlink \hwlink for the three kinds of live link.

% -- Packages ----------------------------------------------------------------
\usepackage[T1]{fontenc}
\usepackage[utf8]{inputenc}
\usepackage{textcomp}
\usepackage[default]{lato}
\usepackage[a4paper,landscape,top=1.25cm,bottom=1.2cm,left=1.4cm,right=1.4cm,%
            headheight=14pt,headsep=10pt,footskip=20pt]{geometry}
\usepackage{xcolor}
\usepackage{colortbl}
\usepackage{array}
\usepackage{pifont}
\usepackage{fancyhdr}
\usepackage{multicol}
\usepackage[most]{tcolorbox}
\usepackage{enumitem}
\usepackage{needspace}
\usepackage{hyperref}

\setlength{\parindent}{0pt}
\setlength{\parskip}{0pt}

% -- The Learner's Book palette (same hex values as decks and homework) -------
\definecolor{cbteal}{HTML}{0087A8}
\definecolor{cbpurple}{HTML}{5C2483}
\definecolor{cborange}{HTML}{C25E12}
\definecolor{cbgreen}{HTML}{4A8B23}
\definecolor{cbred}{HTML}{C8102E}
\definecolor{cbblue}{HTML}{1A5FA8}
\definecolor{rulegray}{HTML}{9AA5AD}
\definecolor{headgray}{HTML}{F2F2F2}

\hypersetup{
  colorlinks=true,
  linkcolor=cbteal,
  urlcolor=cbblue,
  bookmarksopen=true,
  bookmarksnumbered=false,
  pdfstartview={FitH},
  pdfpagelayout=SinglePage,
}

% -- Text-mode substitutes for the three glyphs that would otherwise need maths
\newcommand{\tminus}{\textminus}
\newcommand{\ttimes}{\texttimes}
\newcommand{\sep}{\,\textperiodcentered\,}
\newcommand{\playmark}{\textcolor{cbteal}{\ding{228}}}

% -- The three kinds of live link ---------------------------------------------
% Slides: the deck the class actually saw, on the deployed site.
% Plan:   the one-page teacher plan for the same unit.
% HW:     the homework shelf. Individual packet PDFs are fingerprinted by Vite
%         at build time, so their URLs change on every deploy -- the shelf page
%         is the only stable address for a packet, and it lists them all.
\newcommand{\siteroot}{https://bowenpra.github.io/classroom/\#}
\newcommand{\slidelink}[2]{\playmark~\href{\siteroot/lesson/#1}{Slides \sep #2}}
\newcommand{\planlink}[2]{\href{\siteroot/plan/#1}{Teacher plan \sep #2}}
\newcommand{\hwlink}[1]{\href{\siteroot/homework}{#1}}

% Slides link for the PARENT UPDATE ONLY. Copying a hyperlinked label out of
% the PDF into a Google Classroom post loses the href -- Classroom only
% auto-links plain URL text -- so here the visible text IS the full address,
% not a friendly label. #1 course/unit (the href target); #2 human label
% (printed above the link, for context); #3 the SAME address as #1, typed out
% by hand with \# and \_ escaped, since a macro argument can't be redisplayed
% as literal text without that escaping.
\newcommand{\slidelinkfull}[3]{%
  \needspace{3\baselineskip}%
  \playmark~Slides \sep #2\par\vspace{2pt}%
  \href{\siteroot/lesson/#1}{#3}}

% -- Column widths for the day table -----------------------------------------
% Ratios carried over from the Word version so the two look like one document.
% They sum to 25.1cm, which is \textwidth less tabcolsep and rules.
\newlength{\colSubj}\setlength{\colSubj}{2.5cm}
\newlength{\colLO}  \setlength{\colLO}{6.0cm}
\newlength{\colAct} \setlength{\colAct}{9.4cm}
\newlength{\colEv}  \setlength{\colEv}{7.2cm}

\newcolumntype{L}[1]{>{\raggedright\arraybackslash}p{#1}}

% -- A line inside a table cell ----------------------------------------------
% Cells are lists of short statements, not paragraphs. \li sets one and leaves
% a gap after it, so a cell is written \li{...}\li{...} with no blank lines.
\newcommand{\li}[1]{#1\par\vspace{3.2pt}}

% -- Running head and foot ----------------------------------------------------
\newcommand{\currentweek}{}
\renewcommand{\currentweek}{}
\pagestyle{fancy}
\fancyhf{}
\fancyhead[L]{\footnotesize\color{black!55}Mr Bowen\sep Year 7 Mathematics \& Science}
\fancyhead[R]{\footnotesize\color{black!55}\currentweek}
\fancyfoot[R]{\footnotesize\color{black!55}\thepage}
\fancyfoot[L]{\footnotesize\color{black!40}\hyperref[toc]{Contents}}
\renewcommand{\headrulewidth}{0pt}
\renewcommand{\headrule}{\vspace{2pt}\hbox to\headwidth{\color{cbteal!45}\leaders\hrule height 1.1pt\hfill}}

% -- Page headings ------------------------------------------------------------
% A big coloured title with a thin rule under it. Used by every page type so
% the document reads as one thing.
\newcommand{\planheading}[2]{%
  \noindent{\LARGE\bfseries\color{#1}#2}\par\vspace{4pt}%
  \hbox to \textwidth{\color{#1!35}\leaders\hrule height 1.6pt\hfill}\par\vspace{10pt}}

% The same thing one step down, for a second block on a page that already has a
% \planheading at the top of it.
\newcommand{\planheadsub}[2]{%
  \noindent{\Large\bfseries\color{#1}#2}\par\vspace{3pt}%
  \hbox to \textwidth{\color{#1!35}\leaders\hrule height 1.3pt\hfill}\par\vspace{8pt}}

% -- The day page -------------------------------------------------------------
% One weekday, one page, always. #1 is the label key used by the contents.
\newenvironment{daypage}[3]{%
  \clearpage
  \phantomsection\label{#1}%
  \pdfbookmark[1]{#2 #3}{#1}%
  \planheading{cbteal}{#2 #3}%
  \footnotesize
  \renewcommand{\arraystretch}{1.25}%
  \arrayrulecolor{black!45}%
  \noindent\begin{tabular}{|L{\colSubj}|L{\colLO}|L{\colAct}|L{\colEv}|}
  \hline
  \rowcolor{headgray}
  \centering\textbf{Subject} & \centering\textbf{Learning Objective / Focus} &
  \centering\textbf{Main Activities} & \centering\arraybackslash\textbf{Evidence of Learning} \\
  \hline
}{%
  \end{tabular}\par
}

% \subjectrow{Homeroom}{8:30 -- 8:45}{objective}{activities}{evidence}
\newcommand{\subjectrow}[5]{%
  \textbf{#1}\par\vspace{1pt}{\color{black!55}#2} & #3 & #4 & #5 \\
  \hline}

% -- The week overview page ---------------------------------------------------
\newenvironment{weekoverview}[3]{%
  \clearpage
  \phantomsection\label{#1}%
  \pdfbookmark[0]{#2 \textperiodcentered\ #3}{#1}%
  \planheading{cbpurple}{#2 \sep #3}%
  \small
  \renewcommand{\arraystretch}{1.3}%
  \arrayrulecolor{black!45}%
  \noindent\begin{tabular}{|L{3.1cm}|L{10.7cm}|L{11.3cm}|}
  \hline
  \rowcolor{headgray}
  \centering\textbf{Day} & \centering\textbf{Science \sep 8:45 -- 9:35} &
  \centering\arraybackslash\textbf{Mathematics \sep 10:45 -- 11:35} \\
  \hline
}{%
  \end{tabular}\par\vspace{14pt}%
}

\newcommand{\glancerow}[4]{%
  \hyperref[#1]{\textbf{#2}} & #3 & #4 \\ \hline}

% A block of links under the at-a-glance table. Two columns, because the list
% is short lines and a single column wastes half a landscape page.
\newenvironment{weeklinks}{%
  {\normalsize\bfseries\color{cborange}\ding{228}~Links for this week}\par\vspace{6pt}%
  \small
  \begin{multicols}{2}
}{%
  \end{multicols}
}
\newcommand{\linkline}[1]{#1\par\vspace{4pt}}

% -- The parent announcement --------------------------------------------------
% Deliberately NOT its own page. It follows the at-a-glance table and the links
% on the week page, because it is week-level material and the overview alone
% leaves two thirds of a landscape page empty. A long update simply flows onto
% a second page; the next \daypage clears it either way.
\newcommand{\annstart}[2]{%
  \par\vspace{15pt}%
  \needspace{7\baselineskip}%
  \phantomsection\label{#1}%
  \pdfbookmark[1]{Parent update -- #2}{#1}%
  \planheadsub{cbgreen}{Parent update \sep #2}}

\newenvironment{annbody}{\small\begin{multicols}{2}}{\end{multicols}}

\newcommand{\annsec}[1]{%
  \par\vspace{9pt}%
  {\bfseries\color{cbgreen}\MakeUppercase{#1}}\par\vspace{2pt}%
  \hbox to \columnwidth{\color{cbgreen!35}\leaders\hrule height 1pt\hfill}\par\vspace{6pt}}

\newcommand{\annsub}[1]{%
  \par\vspace{6pt}{\bfseries\color{black!80}#1}\par\vspace{2pt}}

% A note from Mr Bowen to himself, never part of the announcement text.
\newtcolorbox{editorial}[1][Note]{%
  enhanced, breakable, arc=5pt, boxrule=1pt,
  colback=cborange!6, colframe=cborange, coltitle=cborange,
  colbacktitle=white, fonttitle=\bfseries\small,
  attach boxed title to top left={xshift=11pt, yshift=-7pt},
  boxed title style={arc=3pt, boxrule=1pt, left=6pt, right=6pt, top=1pt, bottom=1pt},
  left=10pt, right=10pt, top=12pt, bottom=8pt, title={#1}, fontupper=\small}

% -- The hand-rolled table of contents ----------------------------------------
% Hand-rolled rather than \tableofcontents because this document's structure is
% weeks and days, not sections, and because \tableofcontents inside multicols
% is fragile. Every entry is a live internal link plus a page number.
\newcommand{\tocweek}[2]{%
  \par\vspace{9pt}%
  {\bfseries\color{cbpurple}\hyperref[#1]{#2}}\par\vspace{3pt}}
% The colour must be set OUTSIDE the leader box. xcolor's \color queues an
% \aftergroup\reset@color, so colouring the box from within lands that reset
% between the leaders and the \hfill and TeX stops with "Leaders not followed
% by proper glue" -- once per contents line.
\newcommand{\tocline}[2]{%
  \noindent\hspace{1.1em}\hyperref[#1]{#2}%
  \nobreak\ {\color{black!30}\leaders\hbox to 0.55em{\hss.\hss}\hfill}\ \nobreak
  \hyperref[#1]{\pageref{#1}}\par\vspace{1.5pt}}


\hypersetup{
  pdftitle={Year 7 Mathematics and Science -- Lesson Plan, Quarter 1},
  pdfauthor={Mr Bowen},
  pdfsubject={Cambridge Lower Secondary Stage 7 -- weekly lesson plan},
}

\begin{document}

% =============================================================================
% TITLE
% =============================================================================
\thispagestyle{empty}
\pdfbookmark[0]{Title}{title}
\vspace*{1.6cm}
\begin{center}
  {\Huge\bfseries\color{cbteal}Year 7 Mathematics \& Science}\par\vspace{10pt}
  {\LARGE\bfseries\color{black!70}Lesson Plan \sep Quarter 1, Weeks 1--6}\par\vspace{6pt}
  {\Large\color{black!55}4 August -- 11 September 2026}\par\vspace{22pt}
  \hbox to 12cm{\color{cbteal!40}\leaders\hrule height 1.6pt\hfill}\par\vspace{20pt}
  {\large\color{black!70}Mr Bowen \sep Cambridge Lower Secondary, Stage 7}\par
\end{center}

\vspace{20pt}
\begin{center}
\begin{minipage}{17cm}
\small
{\bfseries\color{cborange}\ding{228}~How to use this document}\par\vspace{6pt}
Every coloured entry is a live link. The contents page jumps to any week or
day; \textcolor{cbblue}{blue} links open the deployed site --- the slides the
class saw, the one-page teacher plan, or the homework shelf.\par\vspace{5pt}
Each week is one overview page (the week at a glance, every link, and the
Friday parent update as it went out) followed by one page per teaching day,
so a single day can be printed and written on. Orange boxes are notes to the
teacher, never part of a post.\par\vspace{5pt}
Timetable: Science 8:45--9:35, Mathematics 10:45--11:35.
\end{minipage}
\end{center}

% =============================================================================
% CONTENTS
% =============================================================================
\clearpage
\phantomsection\label{toc}
\pdfbookmark[0]{Contents}{contents}
\planheading{cbteal}{Contents}

\begin{multicols}{2}
\small

\tocweek{w1}{Week 1 \sep 4--7 August 2026}
\tocline{w1}{Overview and links}
\tocline{w1ann}{Parent update -- Friday 7 August}
\tocline{w1d1}{Tuesday 4 August}
\tocline{w1d2}{Wednesday 5 August}
\tocline{w1d3}{Thursday 6 August}
\tocline{w1d4}{Friday 7 August}

\tocweek{w2}{Week 2 \sep 10--14 August 2026}
\tocline{w2}{Overview and links}
\tocline{w2ann}{Parent update -- Friday 14 August}
\tocline{w2d1}{Monday 10 August}
\tocline{w2d2}{Tuesday 11 August}
\tocline{w2d3}{Wednesday 12 August}
\tocline{w2d4}{Thursday 13 August}
\tocline{w2d5}{Friday 14 August}

\tocweek{w3}{Week 3 \sep 17--21 August 2026}
\tocline{w3}{Overview and links}
\tocline{w3ann}{Parent update -- Friday 21 August}
\tocline{w3d1}{Monday 17 August}
\tocline{w3d2}{Tuesday 18 August}
\tocline{w3d3}{Wednesday 19 August}
\tocline{w3d4}{Thursday 20 August}
\tocline{w3d5}{Friday 21 August}

\tocweek{w4}{Week 4 \sep 24--28 August 2026}
\tocline{w4}{Overview and links}
\tocline{w4ann}{Parent update -- Friday 28 August}
\tocline{w4d1}{Monday 24 August}
\tocline{w4d2}{Tuesday 25 August}
\tocline{w4d3}{Wednesday 26 August}
\tocline{w4d4}{Thursday 27 August}
\tocline{w4d5}{Friday 28 August}

\tocweek{w5}{Week 5 \sep 31 August--4 September 2026}
\tocline{w5}{Overview and links}
\tocline{w5ann}{Parent update -- Friday 4 September}
\tocline{w5d1}{Thursday 3 September}
\tocline{w5d2}{Friday 4 September}

\tocweek{w6}{Week 6 \sep 7--11 September 2026}
\tocline{w6}{Overview and links}
\tocline{w6ann}{Parent update -- Friday 11 September}
\tocline{w6d1}{Monday 7 September}
\tocline{w6d2}{Tuesday 8 September}
\tocline{w6d3}{Wednesday 9 September}
\tocline{w6d4}{Thursday 10 September}
\tocline{w6d5}{Friday 11 September}

\end{multicols}

% #############################################################################
% WEEK 1
% #############################################################################
\clearpage
\renewcommand{\currentweek}{Week 1 \sep 4--7 August}

\begin{weekoverview}{w1}{Week 1}{4--7 August 2026}
\glancerow{w1d1}{Tuesday 4}
  {Day One. Stage 6 topics that return in Stage 7; using a Table of Contents.}
  {Day One. Class expectations, the weekly schedule, setting up equipment.}
\glancerow{w1d2}{Wednesday 5}
  {Unit 1.1 Plant cells, Part 1 --- cell, organelle, and the scale of a cell.}
  {Stage 6 diagnostic --- baseline of what carries over from Year 6.}
\glancerow{w1d3}{Thursday 6}
  {Unit 1.1 Plant cells, Part 2 --- the organelles, and what a microscope does.}
  {Unit 1.1 Adding and subtracting integers.}
\glancerow{w1d4}{Friday 7}
  {Consolidating 1.1 --- Workbook, and the scale-of-a-cell video.}
  {Consolidating 1.1 --- Workbook, then Around the World.}
\end{weekoverview}

\begin{weeklinks}
\linkline{\slidelink{y7-science/U00_1}{Day One: Welcome to Year 7}}
\linkline{\slidelink{y7-science/U01_1}{Science 1.1 Cells}}
\linkline{\planlink{y7-science/U01_1}{Science 1.1}}
\linkline{\slidelink{y7-math/U01_1}{Maths 1.1 Adding \& Subtracting Integers}}
\linkline{\planlink{y7-math/U01_1}{Maths 1.1}}
\linkline{Homework: \hwlink{Packet 1 \sep Integers and Cells} --- out Thursday
  6 August, due Monday 10 August.}
\end{weeklinks}

\annstart{w1ann}{Friday 7 August}
\begin{annbody}

Welcome to our Google Classroom! I'll post a short update here every Friday, so
you always know what we covered that week, what's coming up next, and what
homework is due --- all in the post itself, so there's nothing else to open.

\annsec{Family notebooks}

Every student carries a Family Notebook home. It holds their homework and any
notes from me, so please look through it with your child each week and sign it.
The Friday post and the notebook say the same thing --- just one on the screen
and one on paper.

\annsec{This week}

\annsub{Science --- Unit 1.1: Cells, The Building Blocks of Life}
We learned that every living thing is made of cells, that cells are too small to
see without a microscope, and the names of the main parts inside one.\par
\vspace{3pt}\slidelinkfull{y7-science/U01_1}{Science 1.1 Cells}%
  {https://bowenpra.github.io/classroom/\#/lesson/y7-science/U01\_1}

\annsub{Maths --- Unit 1.1: Adding and Subtracting Integers}
We added and subtracted positive and negative numbers on a number line, and
practised turning an English sentence into a calculation.\par
\vspace{3pt}\slidelinkfull{y7-math/U01_1}{Maths 1.1 Adding \& Subtracting Integers}%
  {https://bowenpra.github.io/classroom/\#/lesson/y7-math/U01\_1}

\vspace{5pt}
These links open the same slides we used in class. Students can click through
them at home to review.

\annsec{Next week}

\annsub{Science}
Animal cells, then specialised cells --- cells with a special shape for a
special job, like a red blood cell or a root hair cell.

\annsub{Maths}
Multiplying and dividing negative numbers (the sign rules), then factors ---
the numbers that divide exactly into another number.

\vspace{5pt}
I'll post these links in next Friday's update.

\annsec{Homework}

Every student was given a paper homework packet on Thursday. It is due
\textbf{Monday}, at the start of class.
\par\vspace{4pt}
\hwlink{Packet 1 \sep Integers and Cells}
\par\vspace{4pt}
The packet has short notes at the top of each section to remind students what
to do, and students can take their notebooks home too if they want extra help
--- our class notes cover everything in the packet. Please make sure it comes
back to school on Monday, and if your child has lost theirs, just let me know
and I'll print another one.

\vspace{8pt}
Thank you,\par
Mr Bowen

\end{annbody}

% -----------------------------------------------------------------------------
\begin{daypage}{w1d1}{Tuesday}{4 August}
\subjectrow{Science}{8:45 -- 9:35}
  {\li{Recognising Stage 6 topics that return in Stage 7.}
   \li{Using a Table of Contents to find a topic in the Learner's Book.}}
  {\li{Model reading the Contents page from the front --- unit, topic, page
      number.}
   \li{Groups scan the Year 7 Contents and pick out topics they already met in
      Year 6.}
   \li{Poster: the topic name, its page number, and what they already know
      about it.}
   \li{Posters shared with the room at the end.}}
  {\li{Poster gives a correct topic with its correct page number.}
   \li{Students can open the book to a named topic without help.}
   \li{Prior knowledge is written in full sentences.}}
\subjectrow{Mathematics}{10:45 -- 11:35}
  {\li{Class expectations, the weekly schedule, and setting up equipment.}
   \li{Students can explain the expectations at home.}}
  {\li{\slidelink{y7-science/U00_1}{Day One, from ``Your Week, Every Week''}}
   \li{Timetable slide --- class reads the four times back: 8:00 arrival, 8:30
      homeroom, 8:45 Science, 10:45 Mathematics.}
   \li{Expectation letter, seven sections, students following on their own
      printed copy.}
   \li{Distribute supplies: binder, folder, three sleeves, Learner's Book.}
   \li{Set up: names on binder and folder; sleeves labelled Maths / Science /
      English; letter at the very front; taped name on the Learner's Book.}}
  {\li{Hold-it-up check --- all six items correct, checked from the front of
      the room.}
   \li{Students can state their two lesson times and where the phone goes on
      arrival.}}
\end{daypage}

% -----------------------------------------------------------------------------
\begin{daypage}{w1d2}{Wednesday}{5 August}
\subjectrow{Science}{8:45 -- 9:35}
  {\li{Unit 1.1 Plant cells, Part 1.}
   \li{Define a cell and an organelle, and describe how small a cell is.}}
  {\li{\slidelink{y7-science/U01_1}{Unit 1.1 Cells, Part 1, up to ``Parts of a
      Plant Cell''}}
   \li{Starter: one full sentence describing how small a cell is.}
   \li{Soda Can Challenge --- guesses first, then the numbers: 0.02\,mm to
      120\,mm, scale factor 6000, Mr Bowen at 10.68\,km against Everest at
      8.85\,km.}
   \li{Copy the definition of ``cell''; Latin \textit{cella}, the prison cell,
      Hooke's cork drawing.}
   \li{Leaf micrograph --- what are the little green circles?}
   \li{Organelle analogy, then the interactive cell diagram.}
   \li{Draw This: the plant cell from p.\,9, all seven labels.}}
  {\li{Scale factor and the 10.68\,km answer worked out correctly on paper.}
   \li{Definitions of ``cell'' and ``organelle'' copied into notebooks.}
   \li{Labelled plant-cell drawing with all seven labels.}}
\subjectrow{Mathematics}{10:45 -- 11:35}
  {\li{Stage 6 diagnostic --- baseline of what students carry over from
      Year 6.}}
  {\li{Papers face down, 40 minutes, silence.}
   \li{Hands up to have a word READ aloud --- never explained.}
   \li{Students who finish early return to yesterday's science posters.}
   \li{Marked tonight so Thursday can be pitched to it.}}
  {\li{Marked diagnostic paper for every student.}
   \li{Gaps recorded by strand; used to set the pace for Quarter 1.}}
\end{daypage}

% -----------------------------------------------------------------------------
\begin{daypage}{w1d3}{Thursday}{6 August}
\subjectrow{Science}{8:45 -- 9:35}
  {\li{Unit 1.1 Plant cells, Part 2.}
   \li{Name the shared and plant-only organelles and their jobs; explain what a
      microscope does and what a model cannot show.}}
  {\li{\slidelink{y7-science/U01_1}{Unit 1.1 Cells, Part 2, from ``Every Cell
      Has These Four''}}
   \li{The four organelles every cell has.}
   \li{Plant-only extras: cell wall, cellulose and sap vacuole; then chloroplast
      and chlorophyll, with real onion and pondweed.}
   \li{Plant and animal cells side by side.}
   \li{Draw This: the labelled microscope from p.\,10; copy ``microscope'' and
      walk the light path.}
   \li{The electron microscope; models and their limitations.}
   \li{Begin Workbook 1.1, pages 2--3.}}
  {\li{13 definitions and 2 labelled drawings in the notebook.}
   \li{Students can say which organelles only plants have.}
   \li{Workbook 1.1 pp.\,2--3 started.}}
\subjectrow{Mathematics}{10:45 -- 11:35}
  {\li{Unit 1.1 Adding and subtracting integers.}
   \li{Move both ways along the number line, including subtracting a negative,
      and turn an English sentence into a calculation.}}
  {\li{\slidelink{y7-math/U01_1}{Maths 1.1 Adding \& Subtracting Integers}}
   \li{Team game: name the negatives against a 2-minute timer, duplicates
      crossed off both lists.}
   \li{Draw This: the number line. Copy ``integer'' and ``inverse''.}
   \li{``Negative five'' against ``minus five'' --- the same dash doing two
      jobs.}
   \li{Jumper widget: pairs say which way the arrow will go before the button
      is pressed.}
   \li{Up-words and down-words, with the sentence patterns.}
   \li{``Subtract 5 from 8'' and ``the difference between \tminus3 and 4'' ---
      hint, wrong answers out loud, then reveal.}
   \li{Word problems on mini-whiteboards: car park, bank account, research
      station, snail.}}
  {\li{8 \tminus\ 5 = 3 given with the reason --- ``from'' marks the number you
      start at.}
   \li{Difference between \tminus3 and 4 is 7, while \tminus3 \tminus\ 4 is
      \tminus7.}
   \li{Calculation written down before the answer on every word problem.}}
\end{daypage}

% -----------------------------------------------------------------------------
\begin{daypage}{w1d4}{Friday}{7 August}
\subjectrow{Science}{8:45 -- 9:35}
  {\li{Consolidating Unit 1.1 independently.}
   \li{Picturing the scale of a cell.}}
  {\li{Workbook 1.1 --- continue and complete, circulating and supporting.}
   \li{Kurzgesagt video: what a cell would be like scaled up enormously.}
   \li{Link the video back to Wednesday's soda-can number.}}
  {\li{Completed Workbook 1.1 section, marked.}
   \li{Students describe one thing the video showed that a drawing cannot.}}
\subjectrow{Mathematics}{10:45 -- 11:35}
  {\li{Consolidating adding and subtracting integers.}
   \li{Recall at speed.}}
  {\li{Workbook 1.1 --- pages taken from the class folders, worked
      independently.}
   \li{``Around the World'' with integer questions if time allows.}}
  {\li{Completed Workbook 1.1 pages.}
   \li{Correct answers at speed in Around the World.}}
\end{daypage}

% #############################################################################
% WEEK 2
% #############################################################################
% \clearpage before the \renewcommand, always: the header is read at shipout,
% so renewing the week name while the previous page is still open stamps the
% NEXT week's name on it.
\clearpage
\renewcommand{\currentweek}{Week 2 \sep 10--14 August}

\begin{weekoverview}{w2}{Week 2}{10--14 August 2026}
\glancerow{w2d1}{Monday 10}
  {Peer-review Homework Packet 1 --- marking a classmate's work.}
  {Unit 1.2 Multiplying \& Dividing Integers, Day 1 --- the four sign rules.}
\glancerow{w2d2}{Tuesday 11}
  {Unit 1.2 Animal Cells --- and how they differ from plant cells.}
  {Unit 1.2, Day 2 --- brackets first, and estimating.}
\glancerow{w2d3}{Wednesday 12}
  {Unit 1.3 Specialised Cells, Day 1 --- the three animal cells.}
  {Unit 1.3 Lowest Common Multiples, Day 1.}
\glancerow{w2d4}{Thursday 13}
  {Unit 1.3, Day 2 --- the two plant cells; the structure--function table.}
  {Unit 1.3, Day 2 --- the ``don't just multiply'' trap; word problems.}
\glancerow{w2d5}{Friday 14}
  {Consolidating 1.2 and 1.3 --- Workbook, independent practice.}
  {Jeopardy review game --- Unit 1.1 to 1.3 in teams.}
\end{weekoverview}

\begin{weeklinks}
\linkline{\slidelink{y7-science/U01_2}{Science 1.2 Animal Cells}}
\linkline{\planlink{y7-science/U01_2}{Science 1.2}}
\linkline{\slidelink{y7-science/U01_3}{Science 1.3 Specialised Cells}}
\linkline{\planlink{y7-science/U01_3}{Science 1.3}}
\linkline{\slidelink{y7-math/U01_2}{Maths 1.2 Multiplying \& Dividing Integers}}
\linkline{\planlink{y7-math/U01_2}{Maths 1.2}}
\linkline{\slidelink{y7-math/U01_3}{Maths 1.3 Lowest Common Multiples}}
\linkline{\planlink{y7-math/U01_3}{Maths 1.3}}
\linkline{\slidelink{games/G02_jeopardy}{Jeopardy --- Friday's review game}}
\linkline{Homework: \hwlink{Packet 2 \sep Signs, Multiples and Cells} --- out
  Thursday 13 August, due Monday 17 August.}
\end{weeklinks}

\annstart{w2ann}{Friday 14 August}

\begin{annbody}

\annsec{This week}

\annsub{Science --- Unit 1.2: Animal Cells, What You Are Made Of}
We named the parts of an animal cell and compared them with plant cells. The
main thing to remember: animal cells have no cell wall and no chloroplasts, so
they're soft and can't make their own food.\par
\vspace{3pt}\slidelinkfull{y7-science/U01_2}{Science 1.2 Animal Cells}%
  {https://bowenpra.github.io/classroom/\#/lesson/y7-science/U01\_2}

\annsub{Science --- Unit 1.3: Specialised Cells, Built for the Job}
We looked at cells with a special shape for a special job --- red blood cells,
nerve cells, root hair cells --- and practised saying why the shape helps:
``It is adapted to \ldots\ because it has \ldots''\par
\vspace{3pt}\slidelinkfull{y7-science/U01_3}{Science 1.3 Specialised Cells}%
  {https://bowenpra.github.io/classroom/\#/lesson/y7-science/U01\_3}

\annsub{Maths --- Unit 1.2: Multiplying \& Dividing Integers}
We learned the four sign rules. The trap we spent time on: two negatives make a
positive when you multiply or divide, but not when you add. \tminus3 +
\tminus4 is \tminus7, not 7.\par
\vspace{3pt}\slidelinkfull{y7-math/U01_2}{Maths 1.2 Multiplying \& Dividing Integers}%
  {https://bowenpra.github.io/classroom/\#/lesson/y7-math/U01\_2}

\annsub{Maths --- Unit 1.3: Lowest Common Multiples}
We listed multiples, found the ones two numbers share, and picked the lowest.
``Common'' in maths means shared, not ordinary.\par
\vspace{3pt}\slidelinkfull{y7-math/U01_3}{Maths 1.3 Lowest Common Multiples}%
  {https://bowenpra.github.io/classroom/\#/lesson/y7-math/U01\_3}

\vspace{5pt}
These links open the same slides we used in class. Students can click through
them at home to review.

\annsec{New: labs start Thursday}

Starting Thursday 20 August, Year 7 will have a science lab every Thursday
morning.

\vspace{4pt}
Our first lab is a Safety and Skills Lab. Before we do any real experiments,
students need to know the rules of the room and how to handle the equipment, so
we'll go through the safety rules, what to do if something spills or breaks, and
how to measure accurately.

\annsec{Next week}

\annsub{Science}
Cells, tissues and organs (1.4). How cells of the same kind group into tissues,
tissues into organs, and organs into systems --- the heart, the lungs, the
stomach. This finishes Unit 1, so we will also review the whole unit before the
Unit 1 check.

\annsub{Maths}
Highest common factors (1.4), then tests for divisibility (1.5). Factors are the
numbers that divide exactly into another number. Then some quick tricks for
checking whether a big number divides by 3, 4, 9 or 11 without doing the
division.

\vspace{5pt}
I'll post these links in next Friday's update.

\annsec{Homework}

Homework Packet 2 went home on Thursday. It is due \textbf{Monday}, at the start
of class, and covers everything above --- the sign rules, multiples and the LCM,
plant and animal cells, and specialised cells.
\par\vspace{4pt}
\hwlink{Packet 2 \sep Signs, Multiples and Cells}
\par\vspace{4pt}
Each section has short notes at the top to remind students what to do. Students
can take their notebooks home for extra help too --- our class notes cover
everything in the packet. If your child has lost theirs, let me know and I'll
print another one.

\annsec{Family notebooks}

Please look through the notebook with your child this week and sign it. The
Friday post and the notebook say the same thing, just one on the screen and one
on paper.

\vspace{8pt}
Thank you,\par
Mr Bowen

\end{annbody}

% -----------------------------------------------------------------------------
\begin{daypage}{w2d1}{Monday}{10 August}
\subjectrow{Science}{8:45 -- 9:35}
  {\li{Peer-review Homework Packet 1 (Cells).}
   \li{Give accurate, kind feedback on another student's work.}}
  {\li{Hand the packets out so each student marks a classmate's copy --- never
      their own.}
   \li{Mark in coloured pencil against the answer key on the board, question by
      question.}
   \li{Tick correct answers; write the correction, not just a cross, where the
      answer is wrong.}
   \li{Total the score, initial the bottom, and return the packet to its owner
      to read.}
   \li{Class tally of the two questions missed most often; reteach those on the
      board.}}
  {\li{Each packet is marked in coloured pencil with corrections, a score and
      the marker's initials.}
   \li{Students can name the two questions the class found hardest.}}
\subjectrow{Mathematics}{10:45 -- 11:35}
  {\li{Unit 1.2 Multiplying \& Dividing Integers, Day 1.}
   \li{Multiply and divide integers; begin the ``two negatives''
      distinction.}}
  {\li{\slidelink{y7-math/U01_2}{Maths 1.2 Multiplying \& Dividing Integers}}
   \li{Starter: work out \tminus3 + \tminus4 and \tminus3 \tminus\
      (\tminus4).}
   \li{Multiplying as repeated adding; the pattern ladder (question only, in
      pairs).}
   \li{Draw This: copy the four sign rules and the \ttimes\ / \textdiv\ sign
      table.}}
  {\li{The four sign rules and the sign table copied into notebooks.}
   \li{Students can explain why a negative \ttimes\ a negative is positive.}}
\end{daypage}

% -----------------------------------------------------------------------------
\begin{daypage}{w2d2}{Tuesday}{11 August}
\subjectrow{Science}{8:45 -- 9:35}
  {\li{Unit 1.2 Animal Cells.}
   \li{Name the parts of an animal cell and how it differs from a plant cell;
      learn the microscope method.}}
  {\li{\slidelink{y7-science/U01_2}{Science 1.2 Animal Cells}}
   \li{Starter: parts of a plant cell beginning with ``c'' (organelles vs
      substances).}
   \li{Parts of an animal cell (Draw This); spot the difference, plant vs
      animal.}
   \li{``Similar is not the same''; Plant-or-Animal photographs; the stain and
      slide-prep method.}}
  {\li{7 written items and 1 labelled drawing in the notebook.}
   \li{Students decide plant vs animal from a photograph, with a reason.}}
\subjectrow{Mathematics}{10:45 -- 11:35}
  {\li{Unit 1.2, Day 2.}
   \li{Consolidate multiplying and dividing; brackets first; estimating.}}
  {\li{Team game: pairs of integers with a product of \tminus24.}
   \li{Work-backwards questions; brackets first, and the addition-inside-brackets
      trap.}
   \li{Estimating by rounding; word problems (the bank fee, the dive, the
      durian).}}
  {\li{Students work the brackets first and estimate to check.}
   \li{Homework set: Workbook 1.1 \& 1.2, pp.\,7--11.}}
\end{daypage}

% -----------------------------------------------------------------------------
\begin{daypage}{w2d3}{Wednesday}{12 August}
\subjectrow{Science}{8:45 -- 9:35}
  {\li{Unit 1.3 Specialised Cells, Day 1.}
   \li{Function, specialised and adapted; the three specialised animal cells.}}
  {\li{\slidelink{y7-science/U01_3}{Science 1.3 Specialised Cells}}
   \li{Starter: finish the four sentences (cell membranes\ldots, a
      nucleus\ldots).}
   \li{Key words: function, specialised, adapted.}
   \li{Red blood cell, neurone and ciliated cell --- drawn diagrams paired with
      real micrographs; Questions 1 and 2.}}
  {\li{Key words copied and the three animal cells described.}
   \li{Students notice that a red blood cell has no nucleus.}}
\subjectrow{Mathematics}{10:45 -- 11:35}
  {\li{Unit 1.3 Lowest Common Multiples, Day 1.}
   \li{Multiple, common multiple, LCM; ``common'' means shared.}}
  {\li{\slidelink{y7-math/U01_3}{Maths 1.3 Lowest Common Multiples}}
   \li{Starter: the multiples of 4 and of 6; the flashing-lights hook.}
   \li{``Common'' means shared, not ordinary; common multiples lead to the
      lowest common multiple.}
   \li{Draw This: Mr Bowen's list method, the LCM of 6 and 9.}}
  {\li{Three definitions and the LCM method copied.}
   \li{Students find the LCM of two numbers by listing.}}
\end{daypage}

% -----------------------------------------------------------------------------
\begin{daypage}{w2d4}{Thursday}{13 August}
\subjectrow{Science}{8:45 -- 9:35}
  {\li{Unit 1.3, Day 2.}
   \li{The two specialised plant cells; structure fits function; build the
      summary table.}}
  {\li{Root hair cell (absorbs water; why it has no chloroplasts) and palisade
      cell (photosynthesis); real plant micrographs.}
   \li{Drill the sentence frame: ``adapted to \rule{1cm}{0.4pt}\ because it has
      \rule{1cm}{0.4pt}.''}
   \li{Draw This: the structure--function table; recall quiz across the five
      cells.}}
  {\li{All five cells entered in the structure--function table.}
   \li{Students use the ``adapted because'' frame in full sentences.}
   \li{Homework: read pp.\,17--21; finish the table; Questions 1--4.}}
\subjectrow{Mathematics}{10:45 -- 11:35}
  {\li{Unit 1.3, Day 2.}
   \li{Consolidate LCM; the ``don't just multiply'' trap; word problems.}}
  {\li{LCM-finder drill; the trap that the LCM of 4 and 8 is 8, not 32.}
   \li{Word problems: the two buses, the barbecue packs, the two taps, the two
      alarms.}}
  {\li{Students know the LCM is usually smaller than the product.}
   \li{Homework set: Exercise 1.3, pp.\,12--13.}}
\end{daypage}

% -----------------------------------------------------------------------------
\begin{daypage}{w2d5}{Friday}{14 August}
\subjectrow{Science}{8:45 -- 9:35}
  {\li{Consolidating Units 1.2 and 1.3 independently.}
   \li{Recall at speed.}}
  {\li{Workbook 1.2 \& 1.3 --- continue and complete, circulating and
      supporting.}}
  {\li{Completed Workbook 1.2 \& 1.3 pages.}}
\subjectrow{Mathematics}{10:45 -- 11:35}
  {\li{Jeopardy review game.}
   \li{Review Unit 1 (integers and multiples, 1.1--1.3) in teams.}}
  {\li{\slidelink{games/G02_jeopardy}{Jeopardy --- Mathematics 1.1--1.3 board}}
   \li{Teams of 3--4. Categories: Adding \& Subtracting, Multiplying \&
      Dividing, Word Problems, Lowest Common Multiple, and ``The English of the
      Question.''}
   \li{A team must show the calculation, not just the answer, to score; point
      values rise down each column.}
   \li{Mr Bowen hosts; a scorekeeper each round; a wagered final question.}
   \li{Close by recapping the two ideas that caught teams out.}}
  {\li{Teams show the calculation, not just the answer.}
   \li{Every student contributes at least one answer.}
   \li{The class can state the sign rules and how to find an LCM.}}
\end{daypage}

% #############################################################################
% WEEK 3
% #############################################################################
\clearpage
\renewcommand{\currentweek}{Week 3 \sep 17--21 August}

\begin{weekoverview}{w3}{Week 3}{17--21 August 2026}
\glancerow{w3d1}{Monday 17}
  {Jeopardy review of Unit 1.1--1.3, before starting 1.4.}
  {Unit 1.4 Highest Common Factors, Day 1.}
\glancerow{w3d2}{Tuesday 18}
  {Unit 1.4 Cells, Tissues and Organs --- the five-level ladder.}
  {Unit 1.4, Day 2 --- the consecutive-numbers investigation.}
\glancerow{w3d3}{Wednesday 19}
  {Unit 1.4, Day 2 --- the organ relay; completing the ladder.}
  {Unit 1.5 Tests for Divisibility, Day 1 --- the nine tests.}
\glancerow{w3d4}{Thursday 20}
  {\textbf{Lab:} safety and skills --- estimating, weighing, and the scalpel.}
  {Unit 1.5, Day 2 --- the shuffle investigation.}
\glancerow{w3d5}{Friday 21}
  {Jeopardy review across Unit 1.1--1.4.}
  {Consolidating 1.5 --- Workbook, then Around the World.}
\end{weekoverview}

\begin{weeklinks}
\linkline{\slidelink{y7-science/U01_4}{Science 1.4 Cells, Tissues and Organs}}
\linkline{\planlink{y7-science/U01_4}{Science 1.4}}
\linkline{\slidelink{y7-math/U01_4}{Maths 1.4 Highest Common Factors}}
\linkline{\planlink{y7-math/U01_4}{Maths 1.4}}
\linkline{\slidelink{y7-math/U01_5}{Maths 1.5 Tests for Divisibility}}
\linkline{\planlink{y7-math/U01_5}{Maths 1.5}}
\linkline{\slidelink{games/G02_jeopardy}{Jeopardy --- Monday and Friday}}
\linkline{Homework: \hwlink{Packet 3 \sep Factors, Tests and Tissues} --- due
  Monday 24 August.}
\linkline{Homework: \hwlink{Extra sheet 3 \sep The Two That Do the Work} ---
  goes out with Packet 3.}
\end{weeklinks}

\annstart{w3ann}{Friday 21 August}

\begin{annbody}

\annsec{This week}

\annsub{Science --- Unit 1.4: Tissues \& Our First Lab}
We wrapped up Unit 1 looking at how cells group into tissues, then ran our first
lab of the year to test our lab safety and measuring skills. Students carefully
measured out the most vital powder in science (coffee), then practised safely
slicing the thinnest possible layers of onion epidermis. We closed out the week
with a round of Review Jeopardy.\par
\vspace{3pt}\slidelinkfull{y7-science/U01_4}{Science 1.4 Cells, Tissues and Organs}%
  {https://bowenpra.github.io/classroom/\#/lesson/y7-science/U01\_4}

\annsub{Maths --- Highest Common Factors \& Tests for Divisibility}
We worked through factors, finding Highest Common Factors, and practised quick
shortcut rules to test if large numbers divide cleanly by 3, 4, 9, or 11.\par
\vspace{3pt}\slidelinkfull{y7-math/U01_4}{Maths 1.4 Highest Common Factors}%
  {https://bowenpra.github.io/classroom/\#/lesson/y7-math/U01\_4}\par
\vspace{3pt}\slidelinkfull{y7-math/U01_5}{Maths 1.5 Tests for Divisibility}%
  {https://bowenpra.github.io/classroom/\#/lesson/y7-math/U01\_5}

\vspace{5pt}
These links open the same slides we used in class. Students can click through
them at home to review.

\annsec{New: the reading challenge}

We're launching our reading challenge! Every completed book gets added to our
classroom reading wall. Encourage your child to pick books they actually want to
read --- building the habit of finishing matters more than the title.

\annsec{Next week}

\annsub{Science}
Unit 1 review, followed by a lab exploring the practical difference between
accuracy and precision.

\annsub{Maths}
Finishing Unit 1 with exponents, then moving directly into our whole-unit
review.

\annsec{Homework}

Homework Packet 3 covers our work with factors, divisibility tests, and cell
tissues. Due \textbf{Monday} at the start of class.
\par\vspace{4pt}
\hwlink{Packet 3 \sep Factors, Tests and Tissues}\par
\vspace{3pt}\hwlink{Extra sheet 3 \sep The Two That Do the Work}
\par\vspace{4pt}
If a student doesn't finish the packet, they may get an additional fundamentals
packet for extra practice.

\vspace{8pt}
Thank you,\par
Mr Bowen

\end{annbody}

% -----------------------------------------------------------------------------
\begin{daypage}{w3d1}{Monday}{17 August}
\subjectrow{Science}{8:45 -- 9:35}
  {\li{Consolidating Unit 1 (1.1--1.3) before starting 1.4.}}
  {\li{\slidelink{games/G02_jeopardy}{Jeopardy --- Science 1.1--1.3 board}}
   \li{Teams of 3--4; categories drawn from cells, specialised cells, and the
      vocabulary from each lesson.}
   \li{Mr Bowen hosts; a scorekeeper each round.}}
  {\li{Teams answer correctly across all three specialised-cell topics.}
   \li{The class can name the two ideas the game caught them out on.}}
\subjectrow{Mathematics}{10:45 -- 11:35}
  {\li{Unit 1.4 Highest Common Factors, Day 1.}
   \li{Find the factors of a number and the highest common factor of two
      numbers by listing; tell a factor apart from a multiple.}}
  {\li{\slidelink{y7-math/U01_4}{Maths 1.4 Highest Common Factors}}
   \li{Starter: pairs multiplying to 12 and to 18, building the two factor
      lists.}
   \li{Hook: 24 pencils and 40 stickers into identical packs --- guesses only,
      answered later in the lesson.}
   \li{Factor vs multiple --- the distinction this lesson turns on.}
   \li{Draw This: HCF of 24 and 80 by listing, Mr Bowen's method.}
   \li{The ``never none'' trap: HCF(8,\,9) is 1, not none; HCF(6,\,18) is 6,
      not 1.}}
  {\li{Common factors and the HCF found correctly by listing.}
   \li{Students can say a factor ``divides in and the list stops'' while a
      multiple ``is landed on and never stops.''}
   \li{Both trap questions answered correctly on mini-whiteboards.}}
\end{daypage}

% -----------------------------------------------------------------------------
\begin{daypage}{w3d2}{Tuesday}{18 August}
\subjectrow{Science}{8:45 -- 9:35}
  {\li{Unit 1.4 Cells, Tissues and Organs.}}
  {\li{\slidelink{y7-science/U01_4}{Science 1.4 Cells, Tissues and Organs}}
   \li{Starter: draw and label a body outline --- brain, heart, stomach,
      intestines, lungs --- then stand and point to each on themselves.}
   \li{Hook: could one ciliated cell alone keep your lungs clean? Paired
      guesses, then the reveal --- millions, joined edge to edge.}
   \li{Copy the definitions of tissue and organ; everyday ``tissue'' vs the
      scientific word.}
   \li{Draw This: the five-level ladder, cell to tissue to organ to organ
      system to organism, one example carried all the way up.}}
  {\li{Tissue and organ defined correctly, including the students' own ``a
      tissue'' vs ``tissue'' sentence pair.}
   \li{Five-level ladder completed with a consistent example.}
   \li{Fingers-up drill: correct answers on the two deliberate traps (a leaf is
      an organ, not a tissue; the palisade layer is a tissue, not an organ).}}
\subjectrow{Mathematics}{10:45 -- 11:35}
  {\li{Unit 1.4, Day 2.}
   \li{Consolidating factors and the highest common factor.}}
  {\li{Team game: HCF-finder pairs against the clock.}
   \li{Workbook 1.4, pp.\,14--15 --- the ribbons, the bookshelf, and the
      consecutive-numbers investigation.}
   \li{Recap the ``never none'' trap before independent work begins.}}
  {\li{HCF found by listing without prompting.}
   \li{Conjecture stated correctly: the HCF of two consecutive numbers is
      always 1.}
   \li{Homework set: Workbook 1.4, pp.\,14--15.}}
\end{daypage}

% -----------------------------------------------------------------------------
\begin{daypage}{w3d3}{Wednesday}{19 August}
\subjectrow{Science}{8:45 -- 9:35}
  {\li{Unit 1.4, Day 2.}
   \li{Consolidating tissues, organs, organ systems and organisms.}}
  {\li{Organ relay in groups of four: one organ each, out loud, in English, no
      repeats.}
   \li{Workbook 1.4 --- complete the five-level ladder; Question 3 run as a hunt
      through the notebook.}
   \li{Recap the breathing and digestive systems, and the organs each shares.}}
  {\li{Organ relay: at least ten organs named correctly as a class.}
   \li{Completed Workbook 1.4 pages.}
   \li{Students can place one example correctly on all five rungs of the
      ladder.}}
\subjectrow{Mathematics}{10:45 -- 11:35}
  {\li{Unit 1.5 Tests for Divisibility, Day 1.}
   \li{Test a number for divisibility by 2, 3, 4, 5, 6, 8, 9, 10 and 11 without
      dividing.}}
  {\li{\slidelink{y7-math/U01_5}{Maths 1.5 Tests for Divisibility}}
   \li{Starter: which of 24, 35, 90, 91 divide by 2? By 5? --- noticing they
      already know two tests.}
   \li{Hook: 3960 against every number from 2 to 11 --- guesses only, paid back
      one test at a time.}
   \li{``6 is a factor of 24'', ``24 is divisible by 6'', ``24 is a multiple
      of 6'' --- one fact, three sentences.}
   \li{The nine tests, copied one at a time, each checked against 3960.}
   \li{Where is 7? --- the number with no easy test, and why.}
   \li{Factor Hunt widget: 1908, 2530, 7128, calling out which test to run.}}
  {\li{All nine tests copied and correctly applied to 3960.}
   \li{Students turn one division fact into all three sentences (factor of /
      divisible by / multiple of).}
   \li{Students can explain why 7 has no test, rather than assuming their notes
      are incomplete.}}
\end{daypage}

% -----------------------------------------------------------------------------
\begin{daypage}{w3d4}{Thursday}{20 August}
\subjectrow{Science}{8:45 -- 9:35}
  {\li{\textbf{Lab: safety and skills.}}
   \li{Practising measuring and cutting safely --- two techniques used all
      year.}}
  {\li{Measuring coffee grounds: eyeball and estimate 10\,g first, then weigh
      accurately on a scale and compare.}
   \li{Scalpel technique: cut the thinnest possible layer of onion epidermis,
      safe hand position, blade always cutting away from fingers.}}
  {\li{Estimated mass compared honestly against the measured mass, in a full
      sentence.}
   \li{Onion epidermis cut thin enough to see through, using correct scalpel
      grip.}
   \li{Every student can state the two safety rules from memory.}}
\subjectrow{Mathematics}{10:45 -- 11:35}
  {\li{Unit 1.5, Day 2.}
   \li{Consolidating the nine divisibility tests.}}
  {\li{Rapid-fire test drill: a number on the board, hands up with the test
      that applies.}
   \li{Workbook 1.4 \& 1.5, independent practice.}
   \li{The shuffle investigation: an arrangement of 8, 3, 6, 1 that fails
      divisibility by 9 --- there isn't one, and why.}}
  {\li{Students say the test out loud, not just the answer --- ``the digits add
      to 18, so it divides by 9.''}
   \li{Completed Workbook 1.4 \& 1.5 pages.}}
\end{daypage}

% -----------------------------------------------------------------------------
\begin{daypage}{w3d5}{Friday}{21 August}
\subjectrow{Science}{8:45 -- 9:35}
  {\li{Consolidating Unit 1 (1.1--1.4) --- recall at speed.}}
  {\li{\slidelink{games/G02_jeopardy}{Jeopardy --- Science 1.1--1.3, with 1.4
      recap questions added}}
   \li{Teams of 3--4; Mr Bowen hosts; a wagered final question.}
   \li{Close by recapping the two ideas that caught teams out.}}
  {\li{Teams show correct answers across cells, specialised cells, and
      tissues/organs.}
   \li{The class can state the five-level ladder from memory.}}
\subjectrow{Mathematics}{10:45 -- 11:35}
  {\li{Consolidating tests for divisibility.}
   \li{Recall at speed.}}
  {\li{Workbook 1.5 --- pages taken from the class folders, worked
      independently.}
   \li{``Around the World'' with divisibility questions if time allows.}}
  {\li{Completed Workbook 1.5 pages.}
   \li{Correct answers at speed in Around the World.}}
\end{daypage}

% #############################################################################
% WEEK 4
% #############################################################################
\clearpage
\renewcommand{\currentweek}{Week 4 \sep 24--28 August}

\begin{weekoverview}{w4}{Week 4}{24--28 August 2026}
\glancerow{w4d1}{Monday 24}
  {Unit 1 test prep --- building the class list of sticking points.}
  {Unit 1.6 Square Roots \& Cube Roots, Day 1.}
\glancerow{w4d2}{Tuesday 25}
  {\textbf{End-of-Unit 1 Test} (1.1--1.4).}
  {Unit 1.6, Day 2 --- trapping a root between two whole numbers.}
\glancerow{w4d3}{Wednesday 26}
  {Independent review practice; the five key words on a poster.}
  {Consolidating 1.6 --- Workbook, then recall at speed.}
\glancerow{w4d4}{Thursday 27}
  {\textbf{Lab:} accuracy vs precision.}
  {\textbf{End-of-Unit 1 Test} (1.1--1.6).}
\glancerow{w4d5}{Friday 28}
  {Plant a seedling; build the classroom planter.}
  {Around the World across the whole of Unit 1.}
\end{weekoverview}

\begin{weeklinks}
\linkline{\slidelink{y7-science/U00_2}{Science: Accuracy and Precision Lab}}
\linkline{\slidelink{y7-math/U01_5}{Maths 1.5 Tests for Divisibility}}
\linkline{\slidelink{y7-math/U01_6}{Maths 1.6 Square Roots \& Cube Roots}}
\linkline{\planlink{y7-math/U01_6}{Maths 1.6}}
\linkline{Homework: \hwlink{Packet 4 \sep Divisibility and Roots} --- due
  Thursday 3 September, at the start of class.}
\end{weeklinks}

\annstart{w4ann}{Friday 28 August}

\begin{annbody}

\annsec{This week}

\annsub{Science --- Unit 1 Review \& Accuracy Precision Lab}
We finished Unit 1 with a short Cambridge assessment covering cells, tissue,
and specializations! Additionally, students did three stations during lab
time to test their ability to measure with both accuracy and precision. Then
on Friday, we made our hanging classroom garden.\par
\vspace{3pt}\slidelinkfull{y7-science/U00_2}{Science: Accuracy and Precision Lab}%
  {https://bowenpra.github.io/classroom/\#/lesson/y7-science/U00\_2}

\annsub{Mathematics --- Unit 1 Review \& Cubes and Squares}
We explored squares, cubes, and roots. Then we ended the week with a Cambridge
assessment.\par
\vspace{3pt}\slidelinkfull{y7-math/U01_5}{Maths 1.5 Tests for Divisibility}%
  {https://bowenpra.github.io/classroom/\#/lesson/y7-math/U01\_5}\par
\vspace{3pt}\slidelinkfull{y7-math/U01_6}{Maths 1.6 Square Roots \& Cube Roots}%
  {https://bowenpra.github.io/classroom/\#/lesson/y7-math/U01\_6}

\vspace{5pt}
These links open the same slides we used in class. Students can click through
them at home to review.

\annsec{The reading challenge update}

Every student in the class has finished at least one book since our challenge
started last week. We are quickly decorating our wall with all the books
finished!

\annsec{Homework Packet 4}

Due \textbf{Thursday}, at the start of class. Please note: any student who
does not complete the packet may receive an additional fundamentals packet
for extra practice. This week's packet is shorter due to the class
participation in the reading challenge.
\par\vspace{4pt}
\hwlink{Packet 4 \sep Divisibility and Roots}

\annsec{Next week (no school Monday, Tuesday, \& Wednesday)}

We have three days holiday this week. Due to the extended time off, I have
made online work for the students to do, reviewing the first mathematics
unit. This is entirely optional as some students have asked for extra
review. Every student has already signed in to this online portal, but
please reach out if you need help.
\par\vspace{4pt}
\href{https://bowenpra.github.io/Dashboard/}{bowenpra.github.io/Dashboard}
\par\vspace{4pt}
Login: Student Name\par
Code: Last 3 Digits of Student Email

\vspace{8pt}
Thank you,\par
Mr Bowen

\end{annbody}

% -----------------------------------------------------------------------------
\begin{daypage}{w4d1}{Monday}{24 August}
\subjectrow{Science}{8:45 -- 9:35}
  {\li{Unit 1 test prep --- getting ready for tomorrow's test.}}
  {\li{Students revisit notebooks and workbook pages from 1.1--1.4.}
   \li{Build a class list of the five key words that come up most (cell,
      organelle, tissue, organ, organism) and one sticking point each.}
   \li{Pairs quiz each other from their own notes.}}
  {\li{A class-generated list of sticking points to target on Tuesday.}
   \li{Students can define at least three of the five key words unprompted.}}
\subjectrow{Mathematics}{10:45 -- 11:35}
  {\li{Unit 1.6 Square Roots \& Cube Roots, Day 1.}
   \li{Square and cube a number, and find a square root and cube root from two
      learned lists.}}
  {\li{\slidelink{y7-math/U01_6}{Maths 1.6 Square Roots \& Cube Roots}}
   \li{Starter: 7 \ttimes\ 7 and 2 \ttimes\ 2 \ttimes\ 2 on whiteboards.}
   \li{Hook: 225 paving stones in a perfect square --- how many along one side?
      Guesses only.}
   \li{Draw This: the twelve square numbers and the six cube numbers, copied in
      full.}
   \li{Squaring and rooting as ``there and back'';
      5\textsuperscript{2} against 5 \ttimes\ 2, the trap that costs marks.}
   \li{Worked example widget: the book's own Worked Example 1.6.}}
  {\li{Twelve squares and six cubes copied correctly.}
   \li{Students can find a square or cube root from the list without a
      calculator.}
   \li{5\textsuperscript{2} read correctly as ``five squared'', not ``five
      times two.''}}
\end{daypage}

% -----------------------------------------------------------------------------
\begin{daypage}{w4d2}{Tuesday}{25 August}
\subjectrow{Science}{8:45 -- 9:35}
  {\li{\textbf{End-of-Unit 1 Test} --- cells, specialised cells, tissues and
      organs (1.1--1.4).}}
  {\li{Papers face down, silence.}
   \li{Hands up to have a word READ aloud --- never explained.}
   \li{Students who finish early check their answers over, then read quietly.}
   \li{Marked and returned before Unit 2 begins.}}
  {\li{Completed end-of-unit test for every student.}
   \li{Marks recorded against 1.1--1.4; gaps flagged before Unit 2.}}
\subjectrow{Mathematics}{10:45 -- 11:35}
  {\li{Unit 1.6, Day 2.}
   \li{Consolidating squares, cubes and roots.}}
  {\li{Team game: name the square or cube root against a timer.}
   \li{Workbook 1.6, independent practice --- trapping a non-whole root between
      two consecutive whole numbers.}
   \li{Recap ``consecutive'' and why it matters for Exercise 1.6.}}
  {\li{Roots trapped correctly between two consecutive whole numbers.}
   \li{Homework set: Exercise 1.6, pp.\,17--19.}}
\end{daypage}

% -----------------------------------------------------------------------------
\begin{daypage}{w4d3}{Wednesday}{26 August}
\subjectrow{Science}{8:45 -- 9:35}
  {\li{Consolidating Unit 1 through independent review practice.}}
  {\li{Workbook Unit 1 review pages, worked independently, circulating and
      supporting.}
   \li{Students who finish early write one full-sentence definition for each of
      the five key words on a poster.}}
  {\li{Completed Unit 1 review workbook pages.}
   \li{Five correct, full-sentence definitions per poster.}}
\subjectrow{Mathematics}{10:45 -- 11:35}
  {\li{Consolidating Unit 1.6.}
   \li{Recall at speed.}}
  {\li{Workbook 1.6 --- pages taken from the class folders, worked
      independently.}
   \li{Rapid-fire recall: squares and cubes against a 2-minute timer.}}
  {\li{Completed Workbook 1.6 pages.}
   \li{Correct answers at speed in the recall drill.}}
\end{daypage}

% -----------------------------------------------------------------------------
\begin{daypage}{w4d4}{Thursday}{27 August}
\subjectrow{Science}{8:45 -- 9:35}
  {\li{\textbf{Lab: accuracy vs precision.}}}
  {\li{Define accuracy (how close to the true value) and precision (how close
      repeated measurements are to each other), with a target diagram.}
   \li{Skills stations: repeat measurements and plot results as
      accurate/precise, accurate/imprecise, inaccurate/precise,
      inaccurate/imprecise.}
   \li{Each group writes one full sentence explaining which of their own
      measurements were accurate, precise, both, or neither.}}
  {\li{Correct classification of results against the target diagram.}
   \li{One full, correct sentence per group distinguishing accuracy from
      precision.}}
\subjectrow{Mathematics}{10:45 -- 11:35}
  {\li{\textbf{End-of-Unit 1 Test} --- integers, factors and multiples,
      divisibility, square and cube roots (1.1--1.6).}}
  {\li{Calculators not allowed.}
   \li{Papers face down, silence.}
   \li{Hands up to have a word READ aloud --- never explained.}
   \li{Marked and returned before Unit 2 begins.}}
  {\li{Completed end-of-unit test for every student.}
   \li{Marks recorded against 1.1--1.6; gaps flagged before Unit 2.}}
\end{daypage}

% -----------------------------------------------------------------------------
\begin{daypage}{w4d5}{Friday}{28 August}
\subjectrow{Science}{8:45 -- 9:35}
  {\li{Plant a seedling; build a classroom planter.}
   \li{Apply what we know about roots and root hair cells; care for a living
      thing.}}
  {\li{Recap: roots and root hair cells absorb water from the soil (links to
      1.3).}
   \li{Each student plants a seedling in soil in their own pot.}
   \li{Build a shared classroom planter; label every pot with a name and
      today's date.}
   \li{Agree a watering rota and write it on the board.}
   \li{Each student writes one full sentence: what a seedling needs to grow.}}
  {\li{Every student has a planted, labelled seedling.}
   \li{A watering rota is agreed and displayed.}
   \li{Each student writes what a plant needs, in a full English sentence.}}
\subjectrow{Mathematics}{10:45 -- 11:35}
  {\li{Consolidating Unit 1 (1.1--1.6).}
   \li{Recall at speed.}}
  {\li{``Around the World'' with questions spanning all of Unit 1 --- integers,
      multiples, factors, divisibility, roots.}
   \li{Mr Bowen hosts; correct and fast beats correct and slow.}}
  {\li{Correct answers at speed in Around the World, across the full range of
      Unit 1 topics.}}
\end{daypage}

% #############################################################################
% WEEK 5
% #############################################################################
\clearpage
\renewcommand{\currentweek}{Week 5 \sep 31 August--4 September}

\begin{weekoverview}{w5}{Week 5}{31 August--4 September 2026}
\textbf{Mon 31 -- Wed 2} & \multicolumn{2}{l|}{No school --- National Day holiday. Optional Unit 1 review on the Dashboard.} \\ \hline
\glancerow{w5d1}{Thursday 3}
  {Unit 1 review --- the five key words, and the questions the test caught.}
  {Unit 1 drills --- integers, factors and multiples, roots, at speed. Packet 4 in.}
\glancerow{w5d2}{Friday 4}
  {Mid-Autumn Festival prep --- booth ideas and crafts for Friday 25 September.}
  {Squish toys from nets of 3D shapes; Unit 1 drills to finish. Packet 5 out.}
\end{weekoverview}

\begin{weeklinks}
\linkline{\slidelink{y7-science/U01_4}{Science 1.4 Cells, Tissues and Organs} (review)}
\linkline{\slidelink{y7-math/U01_6}{Maths 1.6 Square Roots \& Cube Roots} (review)}
\linkline{\href{https://bowenpra.github.io/Dashboard/}{Dashboard \sep optional Unit 1 review online}}
\linkline{Homework: \hwlink{Packet 5 \sep Unit 1, All of It} --- out Friday
  4 September, due Thursday 10 September.}
\linkline{Homework: \hwlink{Extra sheet 5 \sep Bring Down the Next Digit} ---
  goes out with Packet 5.}
\end{weeklinks}

\annstart{w5ann}{Friday 4 September}

\begin{editorial}[Two things to fix next time]
The sent post does not say which period the Mid-Autumn project time used; the
day page puts the booth planning in Science and the nets in Mathematics, which
is a guess. The post also linked the packet by its direct PDF address
(\texttt{assets/hw05-\ldots.pdf}), which changes on every deploy --- the shelf
link below is the one that lasts.
\end{editorial}

\begin{annbody}

\annsub{Science --- Unit 1 Review}
A short week back from the holiday, so we used the extra time for more Unit 1
review.

\annsub{Mathematics --- Unit 1 Drills}
We did more Unit 1 math drills to keep skills sharp after the break.

\annsec{Mid-Autumn Festival prep begins}

We started project time for our upcoming Mid-Autumn Festival on Friday the
25th! Students discussed business ideas and crafts for their booths, and some
made squish toys using net diagrams of 3D shapes.

\annsec{Homework}

Due Thursday at the start of class.
\par\vspace{4pt}
\hwlink{Packet 5 \sep Unit 1, All of It}

\annsec{Next week}

We start Unit 2 in both Math and Science! We'll also have lab next Thursday.

\vspace{8pt}
Thank you,\par
Mr Bowen

\end{annbody}

% -----------------------------------------------------------------------------
\begin{daypage}{w5d1}{Thursday}{3 September}
\subjectrow{Science}{8:45 -- 9:35}
  {\li{Unit 1 review after the holiday.}
   \li{Recall the five key words and the ideas the Unit 1 test caught.}}
  {\li{Starter: define cell, organelle, tissue, organ, organism from memory,
      then check against the notebook.}
   \li{Go over the test questions the class lost most marks on; reteach on
      the board.}
   \li{Pairs quiz each other from Workbook 1.1--1.4; Around the World to
      finish.}}
  {\li{Five key words defined correctly, unprompted.}
   \li{The two weakest test questions answered correctly on whiteboards.}}
\subjectrow{Mathematics}{10:45 -- 11:35}
  {\li{Unit 1 drills after the holiday.}
   \li{Recall at speed across 1.1--1.6.}}
  {\li{Packet 4 collected at the start of class.}
   \li{Rapid-fire rounds on whiteboards: integer sums and products, LCM and
      HCF by listing, divisibility tests, squares and cubes.}
   \li{Reteach any round where more than a third of boards are wrong.}}
  {\li{Packet 4 in from every student; missing packets noted.}
   \li{Correct answers at speed in every round.}}
\end{daypage}

% -----------------------------------------------------------------------------
\begin{daypage}{w5d2}{Friday}{4 September}
\subjectrow{Science}{8:45 -- 9:35}
  {\li{Mid-Autumn Festival prep, session 1.}
   \li{Plan a booth: a product or craft, and who does what.}}
  {\li{Explain the festival day (Friday 25 September) and what a booth
      needs.}
   \li{Groups brainstorm business ideas and crafts; one idea agreed per
      group.}
   \li{Each group lists materials, jobs and a price in full English
      sentences.}}
  {\li{One booth plan per group with a product, a materials list and named
      jobs.}}
\subjectrow{Mathematics}{10:45 -- 11:35}
  {\li{Nets of 3D shapes, applied.}
   \li{Finish the Unit 1 drills; set Packet 5.}}
  {\li{Squish toys: cut, fold and stuff a net of a cube or cuboid; name the
      faces, edges and vertices while building.}
   \li{Remaining Unit 1 drill rounds for anyone not building.}
   \li{Packet 5 and Extra sheet 5 handed out --- due Thursday 10 September.}}
  {\li{A finished squish toy from a correct net.}
   \li{Every student has Packet 5 with the due date written on it.}}
\end{daypage}

% #############################################################################
% WEEK 6
% #############################################################################
\clearpage
\renewcommand{\currentweek}{Week 6 \sep 7--11 September}

\begin{weekoverview}{w6}{Week 6}{7--11 September 2026}
\glancerow{w6d1}{Monday 7}
  {Unit 2.1a Solids, Liquids and Gases --- properties, the two syringes, the
   Word Wall.}
  {Unit 2.1 Constructing Expressions --- a letter for a number you cannot
   count.}
\glancerow{w6d2}{Tuesday 8}
  {Unit 2.1b Particle Theory --- three arrangements; you are the particles.}
  {Unit 2.1, Day 2 --- the four phrases; Workbook 2.1.}
\glancerow{w6d3}{Wednesday 9}
  {Unit 2.2a Changes of State --- the five changes, verb and noun; the particle
   model.}
  {Unit 2.2 Using Expressions and Formulae --- substituting; the hidden times
   sign.}
\glancerow{w6d4}{Thursday 10}
  {Unit 2.2b Measuring, and Heating Water --- meniscus, thermometer, the
   heating curve. Packet 5 in.}
  {Unit 2.2, Day 2 --- formulae with two letters; Workbook 2.2. Packet 6 out.}
\glancerow{w6d5}{Friday 11}
  {Task: The Formula Machine, Part 1 --- temperature scales and the distance to
   a storm.}
  {Task: The Formula Machine, Part 2 (the pour), then Girl Math.}
\end{weekoverview}

\begin{weeklinks}
\linkline{\slidelink{y7-science/U02_1a}{Science 2.1a Solids, Liquids and Gases}}
\linkline{\planlink{y7-science/U02_1a}{Science 2.1a}}
\linkline{\slidelink{y7-science/U02_1b}{Science 2.1b Particle Theory}}
\linkline{\planlink{y7-science/U02_1b}{Science 2.1b}}
\linkline{\slidelink{y7-science/U02_2a}{Science 2.2a Changes of State}}
\linkline{\planlink{y7-science/U02_2a}{Science 2.2a}}
\linkline{\slidelink{y7-science/U02_2b}{Science 2.2b Measuring, and Heating Water}}
\linkline{\planlink{y7-science/U02_2b}{Science 2.2b}}
\linkline{\slidelink{y7-science/U02_model_states}{Science: Heating \& Cooling, a Particle Model} (widget)}
\linkline{\slidelink{games/G01_word-wall?level=solid-liquid-gas}{Word Wall --- solid, liquid, gas} (widget)}
\linkline{\slidelink{y7-math/U02_1}{Maths 2.1 Constructing Expressions}}
\linkline{\planlink{y7-math/U02_1}{Maths 2.1}}
\linkline{\slidelink{y7-math/U02_2}{Maths 2.2 Using Expressions and Formulae}}
\linkline{\planlink{y7-math/U02_2}{Maths 2.2}}
\linkline{\slidelink{y7-math/T01_formula-machine}{Task 1 --- The Formula Machine}}
\linkline{\planlink{y7-math/T01_formula-machine}{Task 1}}
\linkline{\slidelink{y7-math/T02_girl-math}{Task 2 --- Girl Math}}
\linkline{\planlink{y7-math/T02_girl-math}{Task 2}}
\linkline{Homework: \hwlink{Packet 6 \sep Letters and Particles} --- out
  Thursday 10 September, due Monday 14 September.}
\linkline{Homework: \hwlink{Extra sheet 6 \sep One Letter at a Time} --- goes
  out with Packet 6.}
\end{weeklinks}

\annstart{w6ann}{Friday 11 September}

\begin{editorial}[Friday's two periods]
Friday is recorded as the Formula Machine in Science (its formulae are the
thermometer and the storm) and the pour plus Girl Math in Mathematics. If the
split was different, change the two Friday rows; the post below does not
depend on it.
\end{editorial}

\begin{annbody}

\annsec{Science --- Unit 2: Solids, Liquids and Gases}

We started Unit 2 this week. Students learned the properties of solids,
liquids and gases, and the particle theory that explains them: everything is
made of tiny particles, and how they are arranged decides whether something
keeps its shape, pours, or can be squashed. Then we moved on to changes of
state (melting, freezing, evaporating, boiling and condensing) and practised
measuring volume and temperature the way a scientist does.\par
\vspace{3pt}\slidelinkfull{y7-science/U02_1a}{Science 2.1a Solids, Liquids and Gases}%
  {https://bowenpra.github.io/classroom/\#/lesson/y7-science/U02\_1a}\par
\vspace{3pt}\slidelinkfull{y7-science/U02_1b}{Science 2.1b Particle Theory}%
  {https://bowenpra.github.io/classroom/\#/lesson/y7-science/U02\_1b}\par
\vspace{3pt}\slidelinkfull{y7-science/U02_2a}{Science 2.2a Changes of State}%
  {https://bowenpra.github.io/classroom/\#/lesson/y7-science/U02\_2a}\par
\vspace{3pt}\slidelinkfull{y7-science/U02_2b}{Science 2.2b Measuring, and Heating Water}%
  {https://bowenpra.github.io/classroom/\#/lesson/y7-science/U02\_2b}

\vspace{5pt}
New this week: an interactive particle model. Heat or cool ice, mercury,
oxygen or iron and watch the particles change state. Students can try it at
home.\par
\vspace{3pt}\slidelinkfull{y7-science/U02_model_states}{Heating \& Cooling: A Particle Model}%
  {https://bowenpra.github.io/classroom/\#/lesson/y7-science/U02\_model\_states}

\annsec{Maths --- Unit 2: Expressions and Formulae}

We started algebra. Students learned to use a letter for a number we don't
know yet, to write expressions like c \tminus\ 50 and 3n, and to substitute a
number into an expression or a formula. The hardest part is the English:
``3 less than t'' is t \tminus\ 3, not 3 \tminus\ t.\par
\vspace{3pt}\slidelinkfull{y7-math/U02_1}{Maths 2.1 Constructing Expressions}%
  {https://bowenpra.github.io/classroom/\#/lesson/y7-math/U02\_1}\par
\vspace{3pt}\slidelinkfull{y7-math/U02_2}{Maths 2.2 Using Expressions and Formulae}%
  {https://bowenpra.github.io/classroom/\#/lesson/y7-math/U02\_2}

\vspace{5pt}
On Friday we did two tasks. The Formula Machine: running real formulas, such
as converting Celsius to Fahrenheit, finding the distance to a storm from the
gap between the flash and the thunder, and working out the depth of water in a
tank. Girl Math: sharing pieces of rice paper by rules like ``twice as many
as'' and ``3 more than'', solved with algebra against the clock.\par
\vspace{3pt}\slidelinkfull{y7-math/T01_formula-machine}{Task: The Formula Machine}%
  {https://bowenpra.github.io/classroom/\#/lesson/y7-math/T01\_formula-machine}\par
\vspace{3pt}\slidelinkfull{y7-math/T02_girl-math}{Task: Girl Math}%
  {https://bowenpra.github.io/classroom/\#/lesson/y7-math/T02\_girl-math}

\annsec{Online Dashboard (optional)}

The online Dashboard is now updated through Maths 2.2 and Science 2.2. This
is optional, but if your child gets stuck on the homework, Science 2.1 and 2.2
on the Dashboard are the best place to look first.
\par\vspace{4pt}
\href{https://bowenpra.github.io/Dashboard/}{https://bowenpra.github.io/Dashboard/}
\par\vspace{4pt}
Login: student name\par
Code: last 3 digits of the student email

\annsec{Homework --- Packet 6}

Due \textbf{Monday 14 September} at the start of class. It covers this
week's Maths and Science. The last page is a challenge that joins the two
subjects together (volume). It is okay if students cannot finish that page ---
they should try it and show their working.
\par\vspace{4pt}
\hwlink{Packet 6 \sep Letters and Particles}

\annsec{Next week}

\annsub{Science}
Unit 2 continues: using the particle model to explain what happens during a
change of state.

\annsub{Maths}
Unit 2 continues: collecting like terms and simplifying expressions.

\vspace{8pt}
Thank you,\par
Mr Bowen

\end{annbody}

% -----------------------------------------------------------------------------
\begin{daypage}{w6d1}{Monday}{7 September}
\subjectrow{Science}{8:45 -- 9:35}
  {\li{Unit 2.1a Solids, Liquids and Gases.}
   \li{Sort a substance by what it does, not how it looks; list the properties
      of each state; use hypothesis and theory correctly.}}
  {\li{\slidelink{y7-science/U02_1a}{Science 2.1a Solids, Liquids and Gases}}
   \li{Starter: two solids, two liquids, two gases in the notebook.}
   \li{Hook: is sand a liquid? Guesses only; answered after the properties.}
   \li{Copy matter, states of matter, property; one state at a time with a
      copy-down panel each.}
   \li{The two syringes --- predict which plunger moves, then push.}
   \li{Draw This: the four-question table, ruled; Questions 1--7.}
   \li{Page 30: hypothesis, theory, particle; the kitchen-smell question left
      open for tomorrow.}
   \li{Word Wall (solid / liquid / gas puzzle) if time allows.}}
  {\li{Table ruled and ticked correctly for all three states.}
   \li{Students say sand is a solid that pours, and why only a gas
      compresses.}
   \li{Hypothesis and theory defined in the notebook.}}
\subjectrow{Mathematics}{10:45 -- 11:35}
  {\li{Unit 2.1 Constructing Expressions.}
   \li{Use a letter for a number you cannot count; leave c \tminus\ 50
      unfinished; write 3s; read ``less than'' and ``subtract from''
      correctly.}}
  {\li{\slidelink{y7-math/U02_1}{Maths 2.1 Constructing Expressions}}
   \li{Starter: how many particles in one drop of water? (science bridge ---
      nobody can write the number, and that is the point).}
   \li{Name it: let n represent the number; first copy-down panel.}
   \li{The wall: the coffee table where the third row cannot finish --- c
      \tminus\ 50 is the answer.}
   \li{Notation: 3 \ttimes\ s is written 3s.}
   \li{The four phrases: more than, less than, times, subtract from; the order
      flips for ``h less than t''.}
   \li{``Subtract the result from 25'' is 25 \tminus\ 3n.}}
  {\li{Eight copy-down panels in the notebook.}
   \li{``h less than t'' written t \tminus\ h on whiteboards.}
   \li{c \tminus\ 50 accepted as a finished answer.}}
\end{daypage}

% -----------------------------------------------------------------------------
\begin{daypage}{w6d2}{Tuesday}{8 September}
\subjectrow{Science}{8:45 -- 9:35}
  {\li{Unit 2.1b Particle Theory.}
   \li{Describe the particle arrangement in each state and use it to explain
      pouring, compressing and a vacuum.}}
  {\li{\slidelink{y7-science/U02_1b}{Science 2.1b Particle Theory}}
   \li{Starter: three things a gas can do that a solid cannot, books closed.}
   \li{Yesterday's kitchen-smell question, then the ink in water.}
   \li{Copy the theory: all matter is made of particles, arranged
      differently.}
   \li{Three arrangements, one copy-down panel each; everyone stands and
      becomes the particles.}
   \li{Page 33: why water pours and a brick does not; vacuum.}
   \li{The sponge: a solid you can squash --- where the air goes.}
   \li{Draw This: the three-box activity; recap and homework.}}
  {\li{Three particle diagrams drawn correctly with properties beside them.}
   \li{Students explain the sponge without saying a solid compresses.}}
\subjectrow{Mathematics}{10:45 -- 11:35}
  {\li{Unit 2.1, Day 2.}
   \li{Consolidate constructing expressions; the four phrases at speed.}}
  {\li{Team game: an English phrase on the board, the expression on
      whiteboards, against a timer.}
   \li{Workbook 2.1 --- independent practice, circulating.}
   \li{Recap ``less than'' and ``subtract from'' before independent work.}}
  {\li{Expressions written correctly for all four phrase types.}
   \li{Completed Workbook 2.1 pages.}}
\end{daypage}

% -----------------------------------------------------------------------------
\begin{daypage}{w6d3}{Wednesday}{9 September}
\subjectrow{Science}{8:45 -- 9:35}
  {\li{Unit 2.2a Changes of State.}
   \li{Name the five changes as journeys from one state to another; verb and
      noun for each; melting and boiling points of water.}}
  {\li{\slidelink{y7-science/U02_2a}{Science 2.2a Changes of State}}
   \li{Starter: draw the three particle pictures from memory.}
   \li{Hook: an ice cube becomes a puddle, the puddle disappears --- where did
      the water go, twice?}
   \li{The from-to frame; the five changes one at a time, with the particle
      model widget showing each one as it happens.}
   \li{The verb / noun pairs (melt / melting); Draw This.}
   \li{Activity 2.2.1 in pairs with the change cards.}}
  {\li{Five changes named as from-to sentences with the right verb and noun.}
   \li{0 and 100 degrees given as water's melting and boiling points.}
   \li{Students distinguish evaporating from boiling.}}
\subjectrow{Mathematics}{10:45 -- 11:35}
  {\li{Unit 2.2 Using Expressions and Formulae.}
   \li{Substitute a number for a letter; 3n with n = 4 is 12, never 34; keep
      the order of operations; know what a formula is.}}
  {\li{\slidelink{y7-math/U02_2}{Maths 2.2 Using Expressions and Formulae}}
   \li{Starter: the cup from 2.1, now c = 320 --- the row finishes.}
   \li{Substitute and value; the football substitution.}
   \li{The wall: the hidden times sign --- 3n is 3 \ttimes\ n.}
   \li{3x + 2 with x = 4: 14 or 18? Show of hands, then the reason.}
   \li{Formula: an expression has no equals sign, a formula does; A = lw.}
   \li{The minus sign travels with its number --- use brackets.}
   \li{The taxi formula: build one, then break one.}}
  {\li{Seven copy-down panels in the notebook.}
   \li{3n with n = 4 gives 12 on every whiteboard.}
   \li{Students say when a formula stops being true.}}
\end{daypage}

% -----------------------------------------------------------------------------
\begin{daypage}{w6d4}{Thursday}{10 September}
\subjectrow{Science}{8:45 -- 9:35}
  {\li{Unit 2.2b Measuring, and Heating Water.}
   \li{Read a measuring cylinder and a thermometer correctly; plan and run the
      heating-water investigation; plot and describe the heating curve.}}
  {\li{Packet 5 collected at the start of class.}
   \li{\slidelink{y7-science/U02_2b}{Science 2.2b Measuring, and Heating Water}}
   \li{Starter: guess three temperatures in degrees C.}
   \li{Hook: two students read the same cylinder as 50 and 47 --- who is
      right?}
   \li{The meniscus: read the bottom at eye level; the thermometer: read the
      top at eye level. Copy both.}
   \li{Predict, then safety; Draw This: the apparatus, bulb in the water.}
   \li{The investigation, reading every minute; plot temperature against
      time; the particle model at the boiling point.}}
  {\li{Packet 5 in from every student; missing packets noted.}
   \li{Cylinder and thermometer read correctly on the two check questions.}
   \li{A plotted heating curve, and the sentence: it rises, then stays at the
      boiling point.}}
\subjectrow{Mathematics}{10:45 -- 11:35}
  {\li{Unit 2.2, Day 2.}
   \li{Consolidate substituting; formulae with two letters; set Packet 6.}}
  {\li{Team game: substitute against the timer --- 3n, 3n + 2, 25 \tminus\
      3n.}
   \li{Workbook 2.2 --- independent practice, circulating.}
   \li{Recap the hidden times sign and the travelling minus before
      independent work.}
   \li{Packet 6 and Extra sheet 6 handed out --- due Monday 14 September;
      say that Section C is a challenge.}}
  {\li{Completed Workbook 2.2 pages.}
   \li{Every student has Packet 6 with the due date written on it.}}
\end{daypage}

% -----------------------------------------------------------------------------
\begin{daypage}{w6d5}{Friday}{11 September}
\subjectrow{Science}{8:45 -- 9:35}
  {\li{Task: The Formula Machine, Part 1 --- formulae from Science.}
   \li{Run a formula in three lines: write it, put the numbers in, work it
      out.}}
  {\li{\slidelink{y7-math/T01_formula-machine}{Task: The Formula Machine}}
   \li{Starter: the ant walks round the rectangle --- 22.}
   \li{The three-line routine, copied; the middle line is where the marks
      are.}
   \li{Station 1, perimeter; Station 2, area of a triangle --- the height is
      the dashed line.}
   \li{The argument: 35 degrees C and 95 degrees F are the same temperature ---
      link to Thursday's thermometer.}
   \li{Station 4, the storm: flash first, then the bang; d = t divided by 3.}}
  {\li{Every problem worked in three lines, not one.}
   \li{Station 3: 35 degrees C converted to 95 degrees F, with the working.}
   \li{Students say why the flash arrives before the bang.}}
\subjectrow{Mathematics}{10:45 -- 11:35}
  {\li{Task: The Formula Machine, Part 2, then Task: Girl Math.}
   \li{Find the letter that was not given by counting cubes; write
      ``twice as many'' and ``3 more than'' as expressions and solve.}}
  {\li{The pour: water from one tank into another with a different base ---
      deeper, shallower or the same? Real water first, then the widget.}
   \li{\slidelink{y7-math/T02_girl-math}{Task: Girl Math}}
   \li{Twice as many is 2n; 3 more than is n + 3; the four steps.}
   \li{The game: teams, ten rounds, read each round aloud in English before
      the clock runs.}
   \li{Exit question: 34 pieces, Su three times Ana, Tess 4 more than Ana.}}
  {\li{Depth found by counting cubes, not by rearranging.}
   \li{Every share written as an expression before the total is used.}
   \li{Exit question solved by at least half the class.}}
\end{daypage}

\end{document}
